\chapter*{General Conclusion}
\addcontentsline{toc}{chapter}{General Conclusion}

This report documented the design and delivery of \projectname{}, an enterprise partnership and project management platform built during an Agile Scrum internship at Capgemini Engineering Tunisie. The objective was clear from the outset: produce a coherent, production-oriented system that unifies the full lifecycle of the company's partner ecosystem, from a public partnership application through scored onboarding, operational relationship management, project delivery tracking, business intelligence, and AI-assisted decision support, all within a single governed environment.

\section*{Delivered Platform}

The resulting platform covers three surfaces. The public website presents the company's partnership value proposition and collects incoming partnership requests through a guided application wizard. The employee portal provides role-differentiated workspaces for five employee profiles, namely administrator, manager, commercial, analyst, and human-resources, each with the appropriate scope of visibility and action. The partner self-service portal gives each active partner a dedicated space to manage its offers, events, documents, contacts, and project engagements.

Across these surfaces, \projectname{} delivers the following functional pillars:

\begin{itemize}
  \item \textbf{Partner lifecycle management} across four partner categories, namely customer, marketing, supplier, and university, with creation, editing, configurable status workflows, and a full auditable status-change history.
  \item \textbf{Role-based access control} enforced at both the routing layer and the data layer, so that each employee role and each partner account can see and act only on the data it is authorized to access.
  \item \textbf{Partnership request review} combining AI-assisted prescoring, manual employee review, structured accept or reject decisions, and optional email notification to the applicant organization.
  \item \textbf{Partner relationship management} covering contacts, offers, events, meetings, documents, notifications, and satisfaction tracking, all linked to the central partner record.
  \item \textbf{Automated partner scoring} through a five-dimension weighted model covering budget, satisfaction, activity, track record, and strategic fit, with an explainable scoring page that renders a score ring, a radar chart, and per-dimension signals justifying the final recommendation.
  \item \textbf{Project portfolio management} with milestones, task tracking, team allocations, document attachment, and health and progress indicators per project.
  \item \textbf{Business intelligence} backed by a dedicated data-warehouse database, serving dashboard charts that summarize portfolio-wide metrics independently of the transactional store.
  \item \textbf{Agentic AI assistant} offering multi-tool streaming responses, a live plan display, persistent conversation history, rich visual results including charts and interactive tables, exportable executive report cards, and deep links to the relevant application pages.
  \item \textbf{Knowledge retrieval} based on document ingestion, chunking with overlap, embeddings, \sourcecode{pgvector} storage, cosine-distance semantic search, and cited-context assembly that grounds factual answers in real documents.
  \item \textbf{Observability and governance} through structured audit trails, role-aware route protection, grounded AI tools, and evaluation feedback, so that sensitive partnership information remains traceable and controllable.
\end{itemize}

\section*{Agile Scrum Execution}

The internship was organized into short, iterative Agile Scrum sprints, each producing a shippable increment. The first sprint established the project skeleton: authentication, session management, role-based routing, and the application shell. Subsequent sprints layered the core partner data model, the lifecycle workflow, and the public request flow. Later sprints introduced the scoring engine, project management, and the business-intelligence layer. The final sprints assembled the AI harness: the agent engine, the typed tools, the retrieval pipeline, the evaluation framework, and the governance layer.

This cadence kept every intermediate state of the product runnable and reviewable. Scope adjustments stayed traceable to specific sprint boundaries, and the product backlog served as a living record of what was prioritized, deferred, or refined based on discoveries made during implementation.

\section*{Technical Architecture and AI Harness}

\projectname{} is built on Next.js 16 App Router with TypeScript in strict mode, Drizzle ORM over PostgreSQL, Tailwind CSS v4, and the Vercel AI SDK for streaming agent interactions. The codebase follows a feature-driven layout: routing in \texttt{app/}, domain services in \texttt{lib/server/services/}, agent tools and retrieval logic in \texttt{lib/server/agent/}, and shared design-system primitives in \texttt{components/ui/}. This separation keeps every layer independently testable and prevents the AI subsystem from entangling with core business logic.

The AI agent runs with typed, Zod-validated tool schemas, streams token by token to the browser, and enforces grounding at the response layer, so that every partner name, ranking, and metric in an agent reply traces back to a real tool result from the live database. The current tool catalog is intentionally read-oriented: it analyses, retrieves, visualises, and reports without directly mutating core partnership records. The evaluation harness runs agent scenarios reproducibly, so that prompt, retrieval, and tool-selection changes can be compared against stable expectations rather than assessed by subjective inspection.

\section*{Limitations}

\paragraph{Intentionally light HR and employee module.}
The human-resources and employee management module, covering the internal employee directory and the partner-linked student-to-employment pipeline, is scoped as a focused foundation in the current release. This was a deliberate product decision. A dedicated Capgemini human-resources and workforce management sister project is under separate development, and the \projectname{} human-resources module is designed to serve as a consumer of that system rather than a competing implementation. Deep-linking the two platforms, establishing shared data contracts, and consolidating the current screens into a proper integration layer belong to that future integration work. Building a full standalone human-resources system in parallel would have diluted the scope of the partnership management platform that is the primary deliverable of this internship.

\section*{Perspectives}

Several natural extensions follow directly from the delivered work.

\paragraph{Deeper HR integration.}
The most immediate perspective is connecting the human-resources module to Capgemini's dedicated human-resources sister project. This integration would enable bidirectional data flows: \projectname{} would consume workforce records for staffing recommendations and partner-linked pipeline tracking, while publishing relevant partnership engagement events back to the human-resources system. Establishing shared API contracts between the two platforms is the primary technical step, and the current human-resources screens already provide a clear integration surface for that work.

\paragraph{Model sovereignty through fine-tuning.}
Because every agent session is traced, the accumulated traces form a growing, curated dataset of real analytical requests and grounded answers. After client approval, this dataset can be used to fine-tune open-source large language models into on-premise, first-class models tailored to Capgemini's partnership domain, giving the company full control over data residency and inference cost while preserving the grounded, tool-driven behaviour of the current agent.

\paragraph{Learned retrieval fusion.}
The retrieval engine already fuses a dense channel with a BM25 lexical channel through Reciprocal Rank Fusion and a re-ranking pass, as described in Chapter~\ref{chap:rag-scoring-observability}. A natural extension is to learn the fusion and re-ranking weights from click and relevance feedback, and to validate them on a larger labelled evaluation set for production-grade confidence intervals.

\paragraph{Open ecosystem through the Model Context Protocol.}
The agent's tools are already typed and validated, which makes them a natural fit for the Model Context Protocol (MCP), the open standard for connecting agents to tools and data \netref{web:mcp}. Exposing the partner, project, scoring, and retrieval tools as an MCP server would let Capgemini's own assistants and third-party agents reuse the platform's governed capabilities, while consuming external MCP servers would extend the agent beyond the current catalog without changing its harness.

\paragraph{Predictive and multimodal analytics.}
The business-intelligence dashboard aggregates data-warehouse tables into charts. A future iteration could add anomaly detection on KPI time series, scheduled predictive churn scoring, and predictive recruitment that anticipates university hiring needs from historical pipelines. Extending ingestion to multimodal documents, so that layout, tables, and figures inside contracts are retrievable alongside text, would broaden the grounded knowledge base, all surfaced through the agent as additional typed tools that follow the existing registration pattern.

\section*{Closing}

\projectname{} demonstrates that a governed, AI-augmented partnership management platform can be designed and delivered to a coherent, production-oriented state within a single academic internship cycle. The combination of a disciplined full-stack architecture, a principled role-based access model, a measurable retrieval-augmented generation pipeline, and a multi-tool agentic layer produces a system that is functional, auditable, explainable, and extensible.

The Agile Scrum cadence allowed scope to be continuously calibrated against what was buildable and valuable at each sprint boundary rather than what was initially imagined on paper. The result is a product whose every delivered feature is runnable, grounded in real data, and observable through the governance tooling, a solid foundation that Capgemini Engineering Tunisie can carry forward into production integration, human-resources system connectivity, and continued platform evolution.
